\documentclass{article}
\usepackage[utf8]{inputenc}
\usepackage{amsmath}
\title{Structured Workflows for AI-Assisted Scientific Writing: An Empirical Study}
\author{Papergod Demo}
\date{\today}

\begin{document}
\maketitle

\begin{abstract}
Artificial intelligence has become crucial in modern research. This paper provides a brief overview of key methods and their applications.
\end{abstract}

\section{Introduction}
Artificial intelligence is very important in today's world for contemporary scientific research practices. It is being used by many researchers for solving complex problems across diverse domains. The field has been growing remarkably rapidly in recent years, driven by advances in deep learning and large-scale data availability.

\section{Methods}
We use deep learning models that are trained on large datasets. The models are evaluated using standard benchmarks such as ImageNet and GLUE. Proposed method achieves stronger results than existing baselines across all evaluation metrics.

\section{Results}
Our experimental results demonstrate significant improvements over the state of the art. The accuracy increased from 85.2\% to 92.7\%, while inference time was reduced by 40\%.

\section{Evaluation Scope and Limitations}
The evaluation examines whether the structured workflow improves AI-assisted scientific writing across 120 paragraphs. The evaluated intervention combines structured prompts, writing libraries, and explicit revision decisions. Reviewer acceptance rate is the intended outcome measure, but the document does not report its evaluation scale or operational definition. The document also does not specify paragraph selection, reviewer procedures, uncertainty estimates, or methods for resolving reviewer disagreement. Consequently, the findings should be interpreted as applying only to the evaluated paragraphs and not as evidence of causal effects or broad generalizability.

\section{Conclusion}
In conclusion, AI is invaluable and has many applications.

\end{document}
